\makeatletter
\@ifundefined{dimmark}{%
  \newcommand{\dimmark}[2]{\marginpar{\footnotesize\color{teal}\textbf{#1}\,#2}}%
}{}
\makeatother

\chapter[Transmettre, stocker, calculer]{Transmettre, stocker, calculer~: l'émergence des techniques de l'information au regard de l'économie et de l'énergie}\label{ch:history}

Ce chapitre retrace comment les techniques de la transmission, du stockage et du calcul se sont inscrites dans les structures de l'économie~--- comment elles ont pesé sur le capital, le travail, l'organisation des marchés ou le financement des infrastructures, entre autres~--- et à quel coût matériel et énergétique elles l'ont fait. Le parcours est chronologique~--- des tablettes d'argile sumériennes au numérique contemporain~--- mais le geste est analytique~: chaque cas historique est observé au prisme de l'économie et du couple énergie-matière, et les questions qu'il soulève émergent du récit lui-même. Ces observations sont rassemblées en fin de chapitre par réduction typologique. À notre connaissance, aucune des traditions intellectuelles mobilisées~--- économie de la croissance, analyse évolutionniste, histoire des systèmes techniques, économie de l'information, économie de l'énergie\footnote{Les cinq traditions et leurs apports respectifs sont développés dans l'introduction générale de la thèse.}~--- n'a constitué le croisement entre structures économiques, infrastructure informationnelle et substrat énergétique en objet d'étude dans la longue durée. Nous pensons pourtant que c'est à leurs croisements que se logent des éléments de structure qu'aucune, prise isolément, ne peut voir.



% =============================================================================
% PRÉAMBULE MÉTHODOLOGIQUE
% =============================================================================

\subsection*{Méthode et protocole}\label{subsec:grammaire}

Le chapitre traverse une dizaine de cas historiques dans l'ordre chronologique. Le mode dominant est l'induction comparative~: chaque cas est observé pour lui-même, les catégories analytiques émergent de la comparaison entre les cas, et les patterns provisoires qu'un cas fait apparaître sont modifiés ou infirmés par le suivant. C'est ce que \textcite{znaniecki_method_1934} appelle l'induction analytique~: non pas la généralisation à partir d'un échantillon, mais la construction progressive de catégories par confrontation séquentielle du concept au cas. Le regroupement final en dimensions économiques n'intervient qu'en fin de parcours, sous forme d'idéaux-types\footnote{\og Il [l'idéal-type] est obtenu par l'accentuation unilatérale d'un ou de plusieurs points de vue et par l'enchaînement d'une multitude de phénomènes individuels, diffus et discrets, présents tantôt en grand nombre, tantôt en petit nombre, par endroits pas du tout, qui s'ordonnent selon ces points de vue unilatéralement choisis, pour former un tableau de pensée homogène\fg{} \parencite[p.~190--191]{weber_objektivitat_1904}.}.

L'induction analytique au sens de Znaniecki n'est pas une protection contre le biais de confirmation~: elle expose au contraire la construction des catégories à la critique en rendant visible le geste sélectif. Le protocole adopté ici en assume la conséquence~: une catégorie n'est conservée que si elle survit au cas suivant~; un cas qui contredit le pattern n'est jamais écarté comme exception mais traité comme révision de la catégorie. Les trois moments abductifs identifiés plus bas sont précisément des points où cette révision a eu lieu.

L'information, dans ses trois modes~--- la transmission, le stockage et le calcul~--- constitue le niveau d'analyse principal~; selon les cas, la technologie est étudiée en elle-même ou dans le système technique plus vaste où elle s'insère et dont elle tire ses effets.

Mais l'induction seule ne produit pas de résultat théorique nouveau~: elle accumule des régularités. Les moments les plus productifs du parcours sont ceux où un cas \emph{contredit} le pattern établi par les précédents. Lorsque sept pays convergent vers le monopole télégraphique, l'induction identifie un pattern~; lorsque la calculatrice mécanique, technologie comparable, produit un marché fragmenté au lieu d'un monopole, le pattern casse. C'est cette rupture qui déclenche un geste inférentiel distinct~: l'abduction, au sens de Peirce\footnote{Charles Sanders Peirce (1839--1914), fondateur du pragmatisme américain. L'abduction, ou \emph{inference to the best explanation}, est formalisée dans ses \emph{Collected Papers} (1931--1958, vol.~5, §~189). Peirce la distingue de l'induction (généralisation à partir de cas) et de la déduction (conséquence d'une règle). Voir \textcite{douven_abduction_2021} pour une présentation contemporaine.}. Un fait surprenant au regard de ce qui précède engendre une hypothèse qui, si elle était vraie, rendrait le fait non surprenant~; l'hypothèse est ensuite mise à l'épreuve dans le cas suivant. \textcite{timmermans_theory_2012} formalisent cette articulation~: l'induction identifie les régularités, l'abduction génère les hypothèses quand les régularités cassent, et la prédiction qui en résulte est testée déductivement dans le cas d'après. Le chapitre suit ce cycle.

Les moments abductifs du parcours sont les suivants~: le contraste entre la calculatrice (marché fragmenté) et la tabulatrice (monopole à 85,7\,\%), qui engendre l'hypothèse que c'est le mode d'appropriation du capital, et non la technologie, qui détermine la structure de marché~; la mécanisation du \emph{Nautical Almanac}, qui coûte plus cher et emploie plus de monde que le calcul humain, ce qui engendre l'hypothèse que la mécanisation réorganise le travail au lieu de le remplacer~; l'écart de quarante-sept ans entre l'invention du \emph{block system} et son adoption obligatoire, qui engendre l'hypothèse que la technique ne se déploie pas sans un véhicule de financement adapté à la structure des coûts et des bénéfices. Ces hypothèses ne sont pas des conclusions~: elles sont les résultats provisoires que le chapitre~2 devra formaliser et que le chapitre~4 devra tester.

Le parcours met également en évidence une direction causale que le plan du chapitre pourrait occulter. Si chaque cas est présenté comme l'irruption d'une technique dans un contexte économique, le matériau montre aussi, de façon récurrente, que la technique est appelée par une demande de coordination que les moyens existants ne satisfont plus~: la coordination territoriale en temps de guerre pour le télégraphe optique, la sécurité ferroviaire sur voie unique pour le \emph{block system}, la gestion d'un recensement devenu ingérable pour la tabulatrice. Cette bidirectionnalité~--- la technique transforme l'économie, mais l'économie appelle la technique~--- n'est pas traitée séparément dans le chapitre~1~; elle deviendra un objet d'analyse à part entière dans le chapitre~2, où les mécanismes de demande induite devront être formalisés au même titre que les mécanismes d'offre technologique\footnote{Cette posture critique de l'innovation-centrisme et l'attention portée aux conditions d'usage et aux configurations institutionnelles de la diffusion partagent avec un courant historiographique récent une orientation commune. \textcite{edgerton_innovation_1998}, dans le numéro \og Histoire des techniques\fg{} dirigé par Y.~Cohen et D.~Pestre pour les \emph{Annales HSS}, en formule la version qui a fait référence en France~; \textcite{winston_media_1998} en propose la modélisation par le couple \og \emph{supervening necessity}\fg{} / \og \emph{suppression of radical potential}\fg{}~; \textcite{gardey_ecrire_2008} en prolonge la version annaliste en étendant l'\og outillage mental\fg{} de \textcite{febvre_outillage_1938} vers la matérialité des inscriptions, des supports et des gestes. Notre démarche s'en distingue par son ancrage CIRED~--- qui suit en parallèle les grandeurs physiques et monétaires~--- et par son objectif final, qui est de produire des mécanismes testables dans un modèle énergie-économie.}.

Comparer des objets aussi éloignés dans le temps et dans la forme suppose une décision qu'il faut expliciter. Ces dispositifs ne partagent ni le substrat matériel, ni le contexte institutionnel, ni même ce que leurs contemporains appelaient \emph{économie}~--- le mot n'a pas le même sens dans la redistribution des palais sumériens, dans le capitalisme industriel victorien et dans les plateformes numériques actuelles. La comparaison ne tient donc pas à une supposée équivalence économique des objets, mais à un geste assumé~: interroger le matériau historique avec des questions que notre époque a les moyens de formuler, en sachant que ces questions sont elles-mêmes datées. C'est ce que Loraux nommait \emph{anachronisme contrôlé}~--- non un défaut à neutraliser, mais une opération productive qui rend visible, par contraste, ce que les régimes antérieurs ne pouvaient pas voir d'eux-mêmes\footnote{\textcite{loraux_elogeanachronisme_1993}~; pour une discussion contemporaine en histoire économique, \textcite{monnet_perspectives_2025}.}. Le contrôle de l'opération tient à deux disciplines~: que l'anachronisme soit nommé comme tel à chaque mobilisation~--- jamais présenté comme évidence intrinsèque du matériau~--- et qu'il soit borné par les sources, c'est-à-dire qu'il ne fasse dire aux acteurs que ce que leurs effets produisaient, jamais ce que leur conscience aurait dû formuler.

Encore faut-il distinguer ce que l'on rétroprojette. Les concepts économiques que mobilise le chapitre~--- externalité de réseau, coûts de transaction, demande dérivée, frontière firme-marché~--- sont des outils du capitalisme industriel et post-industriel~; appliqués à des dispositifs antérieurs, ils ne décrivent pas ce que leurs acteurs faisaient mais ce que leurs effets \emph{produisaient} au sens que notre vocabulaire permet de formuler. La tripartition qui structure ce chapitre~--- \emph{transmettre, stocker, calculer}~--- est dans une situation analogue, mais d'une autre profondeur~: elle ne nous vient pas d'une économie ancienne dont nous aurions changé le vocabulaire, elle naît avec la cybernétique, à la fin de la période même que ce chapitre traite. Babbage en~1834 sépare le \emph{store} et le \emph{mill} dans l'\emph{Analytical Engine}~; \textcite{vonneumann_first_1945} formalise la séparation entre mémoire, unité arithmétique et entrées-sorties~; \textcite{wiener_cybernetics_1948} définit le contrôle comme communication augmentée d'une rétroaction, laquelle exige stockage de l'état passé et calcul de la correction~; \textcite{kittler_grammophon_1986} pose la tripartition \emph{Speicherung~/ \"Ubertragung~/ Verarbeitung} comme structure médiologique des techniques modernes~; \textcite{hilbert_worlds_2011} sont les premiers à les quantifier comme trois grandeurs mondiales distinctes, dont les rythmes de croissance entre~1986 et~2007 (23, 28 et 58~\% par an respectivement) suggèrent qu'il s'agit bien de fonctions \emph{séparables}\footnote{La distinction prolonge la partition \emph{communication technology~/ information technology} de \textcite{vanderkooij_invention_2015} en explicitant le stockage comme fonction propre. Elle est par ailleurs \emph{orthogonale} à la hiérarchie donnée~$\rightarrow$~information~$\rightarrow$~connaissance~$\rightarrow$~sagesse de \textcite{ackoff_data_1989}, reprise par \textcite{fouquet_digitalisation_2023} pour ses séries longues énergie-information~: chacun de ces niveaux peut être transmis, stocké ou calculé. Sur le plan mathématique, \textcite{bailly_mathematics_2011} et \textcite{longo_computing_2011} fournissent le fondement de la distinction discret/continu qui sous-tend l'argument que la transition analogique/numérique est une différence de nature, non de degré~; \textcite{spreng_interdependency_2015} théorise par ailleurs la substituabilité entre énergie, temps et information comme trois intrants fondamentaux de la production~--- un cadre complémentaire au nôtre, qui décompose non les intrants mais les fonctions du système informationnel.}. La rétroprojection de cette grille sur le réseau Chappe ou sur la mécanographie hollerithienne n'est donc pas neutre~: c'est une lecture que la cybernétique rend possible, et qui dévoile dans le matériau ancien des structurations dont les acteurs ne disposaient pas pour se penser eux-mêmes.

C'est sur la question de l'énergie-matière, cependant, que la rétroprojection devient elle-même un objet de la thèse. La catégorie d'\emph{énergie} comme grandeur unifiée n'existe pas avant l'unification thermodynamique du milieu du XIX\textsuperscript{e}~siècle (Mayer, Joule, Helmholtz)~; son importation dans la pensée économique sous la forme d'un processus productif converti en équivalent énergétique reste marginale, depuis les tentatives de Podolinsky en~1880 jusqu'aux fondateurs de l'économie écologique au XX\textsuperscript{e}~siècle \parencite{vianna_franco_history_2023}. Avant 1845, ce n'est pas que les acteurs ou les économistes \emph{évacuaient} l'énergie~: la catégorie qui aurait permis de la nommer comme telle n'était pas disponible. Le silence des comptes hanovriens du Calenberg ou des budgets du \emph{Nautical Almanac} sur un poste \emph{énergie} n'est pas un oubli mais une impossibilité conceptuelle, et cette impossibilité est elle-même historique. Ce que la thèse identifiera plus loin comme un changement de régime du couplage énergie-information se laisse, à cette aune, formuler autrement~: le moment où la coordination énergie-information cesse d'appartenir au seul historien qui la lit \emph{après coup}, pour entrer dans le calcul des acteurs eux-mêmes. La transversale \textbf{0-TER} n'est donc pas une grille ahistorique appliquée uniformément à toutes les périodes~: c'est une catégorie dont l'apparition même fait partie de la chronologie que ce chapitre cherche à reconstituer. La seconde transversale, la \emph{gouvernance} (\textbf{0-GOV}), est dans un statut différent~: elle désigne l'ensemble des modes de coordination autoritaire qui prescrivent et contraignent le système informationnel~--- l'État souverain, l'État régulateur, l'organisme international, la normalisation technique, la corégulation public-privé, l'autorégulation industrielle. Ces modes traversent toutes les périodes mais sous des formes historiquement contrastées, et l'un des résultats que le chapitre cherche à établir est qu'ils ne se substituent pas mais se déplacent~: la modalité souveraine-westphalienne pour le télégraphe Chappe, la régulation antitrust pour le téléphone Bell et la mécanographie hollerithienne, l'allocation internationale du spectre pour le wireless, la gouvernance multi-stakeholder pour Internet, le retour de la modalité souveraine à la période contemporaine à travers les politiques de Cloud souverain, l'\emph{AI Act} et le contrôle des semi-conducteurs. Nommer cette transversale \emph{gouvernance} plutôt que \emph{souveraineté} signale d'emblée cette pluralité des modes et permet, dans la suite de la thèse, l'articulation avec les leviers politiques modélisables (taxes, subsides, régulations, standards techniques, commande publique) qui structureront la cartographie empirique du chapitre~3 et les scénarios du chapitre~4\footnote{Les promesses qui mobilisent le financement (capital privé) et la commande étatique ont, plus largement, une dimension performative que Jasanoff (2015) théorise sous le nom de \emph{sociotechnical imaginaries}~--- visions collectivement tenues, institutionnellement stabilisées et publiquement performées de futurs désirables, qui orientent le déploiement des technologies. La promesse de dématérialisation qui accompagne aujourd'hui le numérique, comme la promesse de souveraineté territoriale qui accompagnait le télégraphe en~1794, relèvent de ce registre. Nous ne mobilisons pas ce cadre interprétatif dans le corps de l'analyse, dont la finalité est de produire des mécanismes testables dans un modèle énergie-économie, mais nous reconnaissons sa pertinence pour qui voudrait prolonger la lecture sur le versant des imaginaires.}.

Le découpage chronologique suit les changements de substrat technique~--- du signal optique au signal électrique, du calcul mécanique au calcul électronique, de l'infrastructure filaire au réseau numérique. Ces bornes ne sont pas des ruptures nettes~; elles marquent des basculements dans les conditions matérielles de la transmission, du stockage et du calcul. Chaque section correspond à un régime technique distinct~--- cette correspondance est une discipline de lecture, pas un postulat~: elle s'ajustera au fil de la thèse.

La sélection des cas suit le principe de la traversée représentative des substrats techniques, pas celui de l'exhaustivité géographique~: chaque section retient le ou les cas qui font apparaître le plus nettement le basculement technique en cause, et les contre-exemples connus (Suisse, Belgique pour le télégraphe~; calculatrice mécanique américaine pour la mécanographie) sont mobilisés là où ils contredisent le pattern qu'établissent les cas dominants. La régularité que le chapitre identifie n'est donc jamais affirmée sur l'absence de contre-exemples mais sur leur intégration argumentée.

Chaque cas est abordé sous cinq aspects~: \emph{technique} (l'artefact et sa trajectoire d'amélioration), \emph{socio-économique} (les marchés, les acteurs, la régulation), \emph{conceptuel} (les cadres théoriques existants, mobilisés quand le cas les appelle), \emph{matériel} (les flux physiques, les extractions, l'empreinte énergétique), et \emph{conditions d'émergence}~: qui demande le déploiement de la technologie, qui le finance et pour quelle promesse, quel contexte institutionnel le rend possible ou le bloque, et pourquoi le déploiement prend cette forme ici et pas ailleurs\footnote{La promesse qui mobilise le financement peut être commerciale (la rente informationnelle du câble transatlantique), militaire (la coordination territoriale du réseau Chappe), scientifique (la précision des éphémérides au \emph{Nautical Almanac}), ou politique (la souveraineté nationale dans le cas Morse au Congrès). Sur la dimension performative de ces promesses, voir la note précédente sur \emph{sociotechnical imaginaries}.}. Les cas de non-émergence~--- la \emph{Difference Engine} de Babbage, financée mais sans demande~; le télégraphe français, bloqué pendant douze ans par le lobby du sémaphore~--- sont traités avec la même attention que les cas d'émergence réussie, parce qu'ils révèlent les conditions nécessaires par leur absence.

Les observations qui émergent du récit sont étiquetées selon trois niveaux de certitude~: \textbf{*}~(fait documenté indépendamment par des sources distinctes), \textbf{+}~(mécanisme proposé comme explication, compatible avec le matériau mais non testé causalement), \textbf{?}~(piste ouverte, pas encore d'argument construit). Ces marqueurs accompagnent la traversée dans les marges et structurent les bilans de fin de section. Le style du chapitre~--- phrases longues, connecteurs logiques, prose continue~--- porte un risque que ces marqueurs sont conçus pour contenir~: le risque de présenter comme nécessaire ce qui est contingent, ou comme démontré ce qui est proposé. Un pattern observé dans sept pays n'est pas un mécanisme~; un mécanisme compatible avec les données n'est pas un résultat causalement identifié~; et un résultat causalement identifié ne vaut que pour le cas et la période qu'il couvre.

Ces marqueurs prolongent un débat méthodologique sur les rapports entre histoire et économie que \textcite{trannoy_economistes_2025} ont récemment cartographié. \textcite{grenier_causalite_2025} y distingue le paradigme de Laplace, pour lequel la régression causale est possible et l'on peut remonter des effets aux causes, et celui de Cournot, pour lequel l'histoire est ordonnée mais contingente et la régression causale périlleuse~; l'historiographie post-laplacienne reconnaît l'existence de relations causales mais doute de la possibilité de les isoler sur un cas singulier. \textcite{trannoy_causalite_2025} prolonge cette prudence en notant qu'un raisonnement par cas ne peut réfuter qu'un seul type de causalité à la fois, nécessaire ou suffisante, jamais les deux. \textcite{monnet_perspectives_2025} identifie trois types de causalité que la thèse mêle sans toujours les nommer~: la causalité \emph{structurelle}, qui articule un modèle théorique et un choc exogène, et que la traversée mobilise quand un mécanisme économique générique est convoqué~; la causalité \emph{interventionniste}, qui repose sur des quasi-expériences, et sous laquelle se rangent les travaux de \textcite{steinwender_real_2018} et \textcite{juhasz_spinning_2018} mobilisés dans les sections suivantes~; et la causalité \emph{processuelle}, qui reconstruit la chaîne d'actions et où le récit est lui-même le véhicule de l'explication, registre dominant du présent chapitre. Une seconde précaution s'ajoute aux marqueurs et porte non sur la régularité observée mais sur le mécanisme qui la produirait. \textcite{sewell_logics_2005}, cité par \textcite{monnet_perspectives_2025}, rappelle que les structures causales sous-jacentes subissent elles-mêmes des mutations dans le temps~: lorsque le chapitre identifie un pattern qui se reproduit sur cent soixante ans, il n'en suit pas que le mécanisme générateur soit stable. La stabilité est une hypothèse à tester dans les chapitres formels, pas un postulat de la traversée historique.

Le regroupement typologique des questions soulevées par le récit intervient en synthèse de fin de chapitre (\S\ref{sec:reduction})~; sa formalisation comme grille de questionnement relève du chapitre~suivant.

% =============================================================================
% TABLEAUX DÉSACTIVÉS — déplacés vers le chapitre 2 (formalisation des dimensions).
% Présents ici en \iffalse...\fi pour archive ; le préambule du chapitre 1 ne les
% mobilise plus, par cohérence avec la posture inductive (les catégories émergent
% du récit, le regroupement typologique intervient en synthèse de fin de chapitre).
% =============================================================================
\iffalse
\begin{table}[htbp]
\caption{Les dix canaux économiques du couplage ICT-économie~--- présentés ici pour le contexte de lecture~; le déploiement systématique relève du chapitre~2}\label{tab:dimensions}
\centering\small
\begin{tabularx}{\textwidth}{@{}cp{3.2cm}Xp{3.8cm}@{}}
\toprule
 & \textbf{Dimension} & \textbf{Questions soulevées par le récit historique} & \textbf{Exemples de travaux} \\
\midrule
D1 & Propriétés du bien informationnel & Quelles propriétés intrinsèques~--- rivalité, codifiabilité, excluabilité, périssabilité, entre autres~--- déterminent la portée des effets économiques de l'information~? & Samuelson 1954, Arrow 1962, Romer 1990, Shapiro \& Varian 1999, Jones \& Tonetti 2020 \\[4pt]
D2 & Rendements d'échelle et externalités de réseau & Le réseau de communication tend-il vers le monopole~? Le coût de production de l'information diminue-t-il ou augmente-t-il avec l'ambition du service~? & Katz \& Shapiro 1985, Arthur 1989, Tirole 1988 \\[4pt]
D3 & Nature du capital informationnel & La technique informationnelle remplace-t-elle du travail par du capital~? Quel type de capital~--- intangible, physique, hybride~? & Jorgenson 2001, Haskel \& Westlake 2018, Field 1992 \\[4pt]
D4 & Travail, qualification et division cognitive & L'information qualifie-t-elle ou déqualifie-t-elle la main-d'\oe{}uvre~? Crée-t-elle de nouvelles tâches ou en détruit-elle~? & Autor 2003, Acemoglu \& Restrepo 2019, Daston 2018 \\[4pt]
D5 & Demande, tarification et préférences & La baisse du coût de l'information engendre-t-elle un rebond de consommation~? La numérisation transforme-t-elle les préférences~? & Schmookler 1966, Jevons 1865, Allcott et al.\ 2020 \\[4pt]
D6 & Signal de prix et intégration des marchés & L'information synchrone réduit-elle les frictions de marché ou crée-t-elle de nouveaux marchés~? L'effet dépend-il de la nature de ce qui est transmis, stocké ou calculé~? & Hayek 1945, Steinwender 2018, Akerlof 1970 \\[4pt]
D7 & Coûts de transaction et formes organisationnelles & Quand le système de prix ne coordonne pas, quelles formes d'organisation émergent~? La technique informationnelle déplace-t-elle la frontière entre la firme et le marché~? & Coase 1937, Williamson 1985, Chandler 1977 \\[4pt]
D8 & Financement et cycles d'infrastructure & Comment les réseaux d'information sont-ils financés~? La séquence investissement massif, faillite, recapitalisation est-elle récurrente~? & Perez 2002, Gordon 2002, Nonnenmacher 2001 \\[4pt]
D9 & Géographie de la production et commerce & Qui produit les composants, les matériaux et l'énergie nécessaires~? La localisation des infrastructures informationnelles suit-elle la géographie de l'énergie~? & Krugman 1991, Headrick 1988, Kander et al.\ 2013 \\[4pt]
D10 & Trajectoires d'innovation et dépendance de sentier & Les choix techniques passés contraignent-ils les choix futurs~? L'innovation incrémentale crée-t-elle des irréversibilités~? & Nelson \& Winter 1982, Rosenberg 1982, Dosi 1982 \\
\bottomrule
\end{tabularx}
\end{table}

\begin{table}[htbp]
\caption{Deux forces transversales~--- l'énergie-matière et la souveraineté~--- qui traversent plusieurs des canaux précédents sans être elles-mêmes des canaux économiques}\label{tab:transversales}
\centering\small
\begin{tabularx}{\textwidth}{@{}cp{3.2cm}Xp{3.8cm}@{}}
\toprule
 & \textbf{Transversale} & \textbf{Questions soulevées par le récit historique} & \textbf{Exemples de travaux} \\
\midrule
0-TER & Énergie, matière et empreinte physique & Quels flux de matière et d'énergie le système informationnel mobilise-t-il~? Quelles sont les extractions, les infrastructures, les externalités environnementales qui le rendent possible~? Ces flux croissent-ils avec l'ambition fonctionnelle, ou peuvent-ils en être découplés~? & Wrigley 2010, Smil 2017, Pottier 2014, Ensmenger 2018 \\[4pt]
0-GOV & Souveraineté, commande étatique et usage militaire-stratégique & Par qui le système informationnel est-il commandité, financé, contraint~? Quelle est la part de l'État-régulateur, de l'État-militaire, de la géopolitique des infrastructures dans la trajectoire technique~? L'information sert-elle le marché, le pouvoir, ou les deux inextricablement~? & Hughes 1983, Headrick 1991, Edwards 1996, Mazzucato 2013 \\
\bottomrule
\end{tabularx}
\end{table}

\begin{table}[htbp]
\caption{Travaux de référence en économie de l'ICT pour chacune des dix dimensions~--- auteurs ayant spécifiquement étudié les effets économiques des techniques de l'information}\label{tab:dimensions-ict}
\centering\small
\begin{tabularx}{\textwidth}{@{}cp{2.8cm}Xp{4.2cm}@{}}
\toprule
 & \textbf{Dimension} & \textbf{Résultat clé en économie de l'ICT} & \textbf{Travaux spécifiques ICT} \\
\midrule
D1 & Propriétés du bien informationnel & Le bien numérique a un coût marginal nul~; les données sont non rivales mais leur valeur dépend de l'agrégation & Shapiro \& Varian 1999, Jones \& Tonetti 2020, Goldfarb \& Tucker 2019 \\[4pt]
D2 & Rendements et externalités de réseau & Les réseaux de communication convergent vers le monopole par un cycle ouverture-fermeture récurrent & Katz \& Shapiro 1985, Economides 1996, Wu 2010 \\[4pt]
D3 & Nature du capital & L'investissement intangible dépasse le tangible aux É.-U. vers 2000~; le capital informationnel se substitue au capital physique avec un levier de 1:7 à 1:13 & Corrado et al.\ 2005, Haskel \& Westlake 2018, Field 1992 \\[4pt]
D4 & Travail et division cognitive & L'ICT polarise l'emploi par substitution des tâches routinières~; l'automatisation du \emph{switching} téléphonique a pris 60~ans & Autor et al.\ 2003, Acemoglu \& Restrepo 2019, Feigenbaum \& Gross 2024 \\[4pt]
D5 & Demande et préférences & Le surplus du consommateur numérique est sous-estimé par le PIB~; les biens gratuits masquent le coût de production & Brynjolfsson et al.\ 2019, Allcott et al.\ 2020, Varian 2018 \\[4pt]
D6 & Signal de prix et marchés & Le télégraphe, le téléphone mobile et Internet créent du commerce qui n'aurait pas eu lieu sans eux~--- l'information est constitutive, pas lubrifiant & Steinwender 2018, Jensen 2007, Aker 2010 \\[4pt]
D7 & Coûts de transaction et organisations & L'IT exige des transformations organisationnelles complémentaires~; la baisse des coûts de transaction est asymétrique selon les dimensions & Brynjolfsson \& Hitt 2000, Bloom et al.\ 2012, Goldfarb \& Tucker 2019 \\[4pt]
D8 & Financement et cycles & Les vagues d'infrastructure ICT suivent le cycle irruption-frenzy-crash-synergy~; la séquence se reproduit du câble transatlantique à l'IA & Perez 2002, Nonnenmacher 2001, Janeway 2012 \\[4pt]
D9 & Géographie et commerce & La distance pèse encore dans le commerce numérique~; la diffusion de l'Internet est géographiquement inégale & Blum \& Goldfarb 2006, Forman et al.\ 2012, Headrick 1988 \\[4pt]
D10 & Innovation et dépendance de sentier & L'ICT peut être analysée comme GPT~; le cadre GPT-énergie n'existe pas encore & Bresnahan \& Trajtenberg 1995, Lipsey et al.\ 2005, Nuvolari 2020 \\
\bottomrule
\end{tabularx}
\end{table}
\fi


% =============================================================================